\documentclass[aspectratio=169,14pt,usenames,dvipsnames]{beamer}

\usepackage[utf8]{inputenc}
\usepackage{enumitem}
\usepackage{calc}

\usepackage{datetime}
\newcommand\builddate{%
   \ifcase \month%
        \or Janeiro%
        \or Fevereiro%
        \or Março%
        \or Abril%
        \or Maio%
        \or Junho%
        \or Julho%
        \or Agosto%
        \or Setembro%
        \or Outubro%
        \or Novembro%
        \or Dezembro%
    \fi\space\number\year%
}

\newcommand{\loadtheme}[1]{%
    \input{themes/#1}%
}
\newcommand{\presentationlanguage}[1]{%
    \usepackage[#1]{babel}%
}

\newcommand{\usecodingsamples}[1]{%
    \usepackage{listings}%
    \input{listings/#1}%
}

% Configura a apresentação para ser executada em tela cheia.
\newcommand{\setfullscreen}{\hypersetup{pdfpagemode=FullScreen}}

% Hide beamer navigation simbols
\beamertemplatenavigationsymbolsempty

%
% Standard frames
%

% coverframe
\newcommand{\coverframe}{%
    \begin{frame} %
        \titlepage %
    \end{frame} %
}

% finalframe{email}
\newcommand{\finalframe}[2][Thank you!]{%
    \begin{frame}%
        \begin{flushright}%
            \huge \textbf{#1}%
            \vfill%
            \large \textbf{#2}%
        \end{flushright}%
    \end{frame}%
}

\newcommand{\bigtitle}[1]{%
    \begin{frame}%
        \begin{center}%
            \Huge {#1}%
        \end{center}%
    \end{frame}%
}



\usecodingsamples{python}

%\loadtheme{apple_keynote_black}
\loadtheme{invaders}

\title{Desenvolvendo Jogos com pygame}
\subtitle{}
\author{Rafael Guterres Jeffman}
\institute{}
\date{2019}

\begin{document}

%01
\coverframe

%02
\bigtitle{Por que \\ jogos?}

%02.a
\begin{frame}
    \frametitle{Desenvolver Jogos}

    \begin{columns}
        \begin{column}{0.5\textwidth}
            \begin{itemize}
                \item É divertido.
                \item Tu sempre quis fazer.
                \item Foi a primeira coisa que tu fez com algo que parecia
                um computador.
            \end{itemize}
        \end{column}
        \begin{column}{0.5\textwidth}
            \begin{itemize}
                \item Não precisa ser difícil.
                \item Não é fácil.
                \item Tu quer mostrar pra todo mundo que tu consegue
                desenvolver um jogo.
            \end{itemize}
        \end{column}
    \end{columns}
    \begin{center}
        \item {\Large \textbf{É muito divertido!}}
    \end{center}
\end{frame}

%03
\bigtitle{Por que \\ Python?}

%03.a
\begin{frame}
    \frametitle{Por que Python?}

    \begin{itemize}
        \item É divertido.
        \item Permite que a preocupação seja o problema.
        \item Faz com que tu aprenda uma linguagem que está sendo muito
        utilizada.
    \end{itemize}
\end{frame}

%04
\bigtitle{Por que \\ pygame?}

%04.a
\begin{frame}
    \frametitle{pygame}

    \begin{itemize}
        \item É multi-plataforma (SDL).
        \item Retira as parada chata da programação de jogos.
        \item Não é um engine de jogos, e queremos programar.
        \item É uma biblioteca de componentes.
    \end{itemize}
\end{frame}

%05
\bigimage{images/moodboard.png}

%06
\begin{frame}
    \frametitle{Every saga has a beginning!}
    \vfill
    \begin{center}
    \large \textbf{Durante um teste de rotina, a nave Genesis é trasportada
    através de um \textit{wormhole} para o quadrante \textit{gamma} da galáxia,
    e precisa sobreviver à \textit{Guerra do Infinito}.}
    \vfill
    \large\textbf{O que era só um dia de testes virou\\uma luta pela sobrevivência.}
    \end{center}
    \vfill
\end{frame}

%06.b
\begin{frame}
    \frametitle{Na vida nada se cria...}

    \begin{itemize}
        \item Gradius/Nemesis
        \item Farscape
        \item Star Trek: Voyager
        \item Um filme muito, muito ruim...
    \end{itemize}
\end{frame}


%07
\begin{frame}[fragile]
    \frametitle{Hello World}

    \begin{python}
        import pygame
        pygame.init()
        # inicia tela
        screen = pygame.display.set_mode((320,200))
        pygame.display.set_caption("Hello World!")
        # loop principal
        running = True
        while running:
            # trata eventos
            for event in pygame.event.get():
                if event.type == pygame.QUIT:
                    running = False
            # atualiza objetos
            # desenha objetos
            # pygame usa double buffer!
            pygame.display.update()
    \end{python}
\end{frame}

%08
\begin{frame}[fragile]
    \frametitle{A janela da aplicação}

    \begin{itemize}
        \item Aplicações pygame podem usar o modo janela ou \textit{fullscreen}.
        \item No modo \textit{fullscreen}, o tamanho da janela é a sua resolução.
    \end{itemize}
    \begin{python}
    width, height = size = (800, 600)
    flags = pygame.FULLSCREEN | pygame.HWSURFACE | pygame.DOUBLEBUF
    screen = pygame.display.set_mode(size, flags)
    \end{python}
\end{frame}

%09
\begin{frame}[fragile]
    \frametitle{Desenhando na tela}
    \begin{itemize}
        \item A estrutura criada pela função \texttt{display.set\_mode} é uma
        superfície, que utilizamos para desenhar na tela.
        \item Esta estrutura pode ser utilizada com o módulo \texttt{pygame.draw}.
    \end{itemize}
    \vfill
    \begin{python}
        pygame.draw.circle(screen, red, (100,100), 50)
        pygame.draw.polygon(screen, white, point_list)
        pygame.draw.rect(screen, white, (x, y, rect_w, rect_h))
        screen.blit(image, (x, y))
    \end{python}
\end{frame}

%10
\begin{frame}
    \frametitle{Um campo de estrelas}

    \begin{itemize}
        \item Todo jogo de plataforma que se preze usa o efeito de \textit{paralaxe}.
        \item Um campo de estrelas com três planos pode ser criado com círculos
        que se movem com velocidades diferentes.
        \item \textit{List comprehensions} e \textit{tuplas} são muito úteis para isso.
    \end{itemize}
\end{frame}

%10.a
\begin{frame}[fragile]
    \frametitle{Um campo de estrelas - Criação}

    \begin{python}
    def create_star(x):
        y = randint(0, height)
        speed = choice([1, 2, 3])
        magnitude = choice([1, 2, 3])
        color = (coice(100, 200, 250),) * 3
        return (x, y, speed, magnitude, color)

    stars = [create_star(randint(0, width)) for star in range(count)]
    \end{python}
\end{frame}

%10.b
\begin{frame}[fragile]
    \frametitle{Um campo de estrelas - Movimentação}

    \begin{python}
    stars = [[x - speed, y, speed, mag, color]
             if x - speed > 0
             else create_star(width)
             for x, y, speed, mag, color in stars]
    \end{python}
\end{frame}

%10.c
\bigimage{images/starfield.png}

%11
\begin{frame}
    \frametitle{Sprites}

    \begin{itemize}
        \item Sprites são imagens 2D, mas nos jogos, eles tem movimento.
        \item O uso de sprites facilita a definição dos objetos móveis
        do jogo.
        \item Sprites, normalmente, tem suporte a transparência.
    \end{itemize}
\end{frame}

%11.a
\bigimage{images/f18-big.png}

%12
\begin{frame}
    \frametitle{Sprites com animação}

    \begin{itemize}
        \item pygame suporta imagens GIF, mas sem animação.
        \item pygame oferece diversos plugins que podem ser utilizados.
        \item Obviamente, existe um plugin para GIF animado.
        \item Mais óbvio ainda... ele estava abandonado e não funcionava...
    \end{itemize}
\end{frame}

%12.a
\begin{frame}
    \begin{columns}
        \begin{column}{0.3\textwidth}
            \includegraphics[width=3cm]{images/out0000.png}
            \vfill
            \includegraphics[width=3cm]{images/out0003.png}
            \vfill
            \includegraphics[width=3cm]{images/out0006.png}
            \vfill
            \includegraphics[width=3cm]{images/out0009.png}
        \end{column}
        \begin{column}{0.3\textwidth}
            \includegraphics[width=3cm]{images/out0001.png}
            \vfill
            \includegraphics[width=3cm]{images/out0004.png}
            \vfill
            \includegraphics[width=3cm]{images/out0007.png}
            \vfill
            \includegraphics[width=3cm]{images/out0010.png}
        \end{column}
        \begin{column}{0.3\textwidth}
            \includegraphics[width=3cm]{images/out0002.png}
            \vfill
            \includegraphics[width=3cm]{images/out0005.png}
            \vfill
            \includegraphics[width=3cm]{images/out0008.png}
            \vfill
            \includegraphics[width=3cm]{images/out0011.png}
        \end{column}
    \end{columns}
\end{frame}

%13
\begin{frame}
    \frametitle{Tratamento de Eventos}

    \begin{itemize}
        \item pygame oferece um sistema de eventos por \textit{polling}.
        \item Para criar um \textit{engine} com um \textit{loop} genérico,
        é preciso permitir que o código cliente seja chamado de volta.
        \item Um mecanismo desses permite que funções cliente sejam chamadas
        para eventos do pygame.
        \item \textit{E para felicidade geral da nação...\\ funções são} cidadãs de primeira ordem\textit{!}
    \end{itemize}
\end{frame}

%13.b
\begin{frame}[fragile]
    \frametitle{O loop de eventos}

    \begin{python}
# loop genérico, em Game.run()
while self.running:
    # handle events
    for event in pygame.event.get():
        handle_event(event)
    \end{python}
\end{frame}

%13.c
\begin{frame}[fragile]
    \frametitle{Respondendo a eventos de teclado}

    \begin{python}
# my code
def move(event):
    """Move player with directional keys."""
    keys = pygame.key.get_pressed()
    dx, dy = 0, 0
    dy = -1 if keys[pygame.K_UP] else 0
    dy = dy + 1 if keys[pygame.K_DOWN] else dy
    dx = -1 if keys[pygame.K_LEFT] else 0
    dx = dx + 1 if keys[pygame.K_RIGHT] else dx
    player.move = (dx * config.speed, dy * config.speed)

# Configuring the game object
game.on_key((pygame.K_UP, pygame.K_DOWN,
             pygame.K_LEFT, pygame.K_RIGHT), move)
    \end{python}
\end{frame}

%14
\bigtitle{E o que mais \\ falta fazer?}

%14.a
\begin{frame}
    \frametitle{Um jogo tem tanta coisa...}

    \begin{itemize}
        \item Tratamento de colisões.
        \item Comportamento de NPCs.
        \item Trocas de fazes.
        \item Cenários.
    \end{itemize}
\end{frame}

% 15
\bigtitle{Sem audio?}

%15.b
\begin{frame}
    \frametitle{pygame Mixer!}

    \begin{itemize}
        \item pygame tem um mixer que, sem configuração, suporta 8 canais de audio.
        \item Suporte a loops de áudio já embutido.
        \item Suporte a diversos formatos de áudio.
        \item Ogg Vorbis é a melhor opção.
        \item E a internet está cheia de loops \textit{royalty free}...
    \end{itemize}
\end{frame}

%16
\begin{frame}
    \frametitle{Quão produtivo é o pygame?}

    \vfill
    \begin{center}
        Você tem 24h para...
        \begin{itemize}
            \item Aprender a usar o pygame.
            \item Escrever uma demonstração.
            \item Criar os \textit{slides} de uma palestra.
            \item Apresentar a palastra.
        \end{itemize}
    \end{center}
    \vfill
\end{frame}

%16.a
\begin{frame}
    \frametitle{Quão produtivo é o pygame?}

    \vfill
    \begin{center}
        E ainda por cima...
        \begin{itemize}
            \item Ministrar três horas de aula.
            \item Dirigir de Porto Alegre a Pelotas (4h).
            \vfill
            \item Passear com os cachorros!
        \end{itemize}
    \end{center}
    \vfill
\end{frame}

% 16.b
\bigtitle{\color{magenta} {\Huge \textbf{E DEU CERTO!}}}

% 17
\bigtitle{E agora? \\ Pra onde vou?}

% 17.a
\bigtitle{\url{https://python.org} \vfill \url{https://pygame.org}}

%18
\begin{frame}
    \frametitle{Próximos Passos}

    \begin{itemize}
        \item Desenvolver um engine para criação de jogos 2D!
        \item Para ensinar programação orientada a objetos com Python.
        \item Para ensinar \textit{design} de jogos.
        \item Para ensinar criação de roteiros de jogos.
    \end{itemize}
\end{frame}

%18.a
\bigtitle{\Large E para criar jogos, né?}

%18.b
\bigtitle{\large Na verdade... \\ Esse nunca foi o objetivo...}

%19
\bigtitle{\Large Mas nem uma demonstração?}

%19.b
\bigtitle{\large \url{https://github.com/rafasgj/genesis}}

%20
\finalframe[Muito Obrigado!]{}

\end{document}
